\IfFileExists{stacks-project.cls}{%
\documentclass{stacks-project}
}{%
\documentclass{amsart}
}
\usepackage{amssymb}

% For dealing with references we use the comment environment
\usepackage{verbatim}
\newenvironment{reference}{\comment}{\endcomment}
%\newenvironment{reference}{}{}
\newenvironment{slogan}{\comment}{\endcomment}
\newenvironment{history}{\comment}{\endcomment}

% For commutative diagrams we use Xy-pic
\usepackage[all]{xy}

% We use 2cell for 2-commutative diagrams.
\xyoption{2cell}
\UseAllTwocells

% We use multicol for the list of chapters between chapters
\usepackage{multicol}

% This is generally recommended for better output
\usepackage{lmodern}
\usepackage[T1]{fontenc}
\usepackage{textalpha}

% For cross-file-references
% We prefer the snapshot tmp/xref/#2.aux,
% which latexmkrc replaces atomically, so
% that documents compiled simultaneously
% never read a half-written .aux file.
% The last argument keeps links to #2.pdf.
\usepackage{xr-hyper}
\newcommand{\maybeexternaldocument}[2]{%
  \IfFileExists{#2.aux}{%
    \externaldocument[#1]{#2}%
  }{%
    \PackageWarning{stack}{Missing #2.aux; external references with prefix #1 unavailable}%
  }%
}
\let\externaldocumentorig\externaldocument
\renewcommand{\externaldocument}[2][]{%
  \IfFileExists{tmp/xref/#2.aux}{%
    \externaldocumentorig[#1]{tmp/xref/#2}[#2.pdf]%
  }{%
  \IfFileExists{#2.aux}{%
    \externaldocumentorig[#1]{#2}%
  }{%
    \PackageWarning{stack}{Missing #2.aux; external references with prefix #1 unavailable}%
  }%
  }%
}

% Package for hypertext links:
\usepackage{hyperref}

% For any local file, say "hello.tex" you want to link to please
% use \externaldocument[hello-]{hello}
\externaldocument[introduction-]{introduction}
\externaldocument[conventions-]{conventions}
\externaldocument[sets-]{sets}
\externaldocument[categories-]{categories}
\externaldocument[topology-]{topology}
\externaldocument[sheaves-]{sheaves}
\externaldocument[sites-]{sites}
\externaldocument[stacks-]{stacks}
\externaldocument[fields-]{fields}
\externaldocument[algebra-]{algebra}
\externaldocument[brauer-]{brauer}
\externaldocument[homology-]{homology}
\externaldocument[derived-]{derived}
\externaldocument[simplicial-]{simplicial}
\externaldocument[more-algebra-]{more-algebra}
\externaldocument[smoothing-]{smoothing}
\externaldocument[modules-]{modules}
\externaldocument[sites-modules-]{sites-modules}
\externaldocument[injectives-]{injectives}
\externaldocument[cohomology-]{cohomology}
\externaldocument[sites-cohomology-]{sites-cohomology}
\externaldocument[dga-]{dga}
\externaldocument[dpa-]{dpa}
\externaldocument[sdga-]{sdga}
\externaldocument[hypercovering-]{hypercovering}
\externaldocument[schemes-]{schemes}
\externaldocument[constructions-]{constructions}
\externaldocument[properties-]{properties}
\externaldocument[morphisms-]{morphisms}
\externaldocument[coherent-]{coherent}
\externaldocument[divisors-]{divisors}
\externaldocument[limits-]{limits}
\externaldocument[varieties-]{varieties}
\externaldocument[topologies-]{topologies}
\externaldocument[descent-]{descent}
\externaldocument[perfect-]{perfect}
\externaldocument[more-morphisms-]{more-morphisms}
\externaldocument[flat-]{flat}
\externaldocument[groupoids-]{groupoids}
\externaldocument[more-groupoids-]{more-groupoids}
\externaldocument[etale-]{etale}
\externaldocument[chow-]{chow}
\externaldocument[intersection-]{intersection}
\externaldocument[pic-]{pic}
\externaldocument[weil-]{weil}
\externaldocument[adequate-]{adequate}
\externaldocument[dualizing-]{dualizing}
\externaldocument[duality-]{duality}
\externaldocument[discriminant-]{discriminant}
\externaldocument[derham-]{derham}
\externaldocument[local-cohomology-]{local-cohomology}
\externaldocument[algebraization-]{algebraization}
\externaldocument[curves-]{curves}
\externaldocument[resolve-]{resolve}
\externaldocument[models-]{models}
\externaldocument[functors-]{functors}
\externaldocument[equiv-]{equiv}
\externaldocument[pione-]{pione}
\externaldocument[etale-cohomology-]{etale-cohomology}
\externaldocument[proetale-]{proetale}
\externaldocument[relative-cycles-]{relative-cycles}
\externaldocument[more-etale-]{more-etale}
\externaldocument[trace-]{trace}
\externaldocument[crystalline-]{crystalline}
\externaldocument[spaces-]{spaces}
\externaldocument[spaces-properties-]{spaces-properties}
\externaldocument[spaces-morphisms-]{spaces-morphisms}
\externaldocument[decent-spaces-]{decent-spaces}
\externaldocument[spaces-cohomology-]{spaces-cohomology}
\externaldocument[spaces-limits-]{spaces-limits}
\externaldocument[spaces-divisors-]{spaces-divisors}
\externaldocument[spaces-over-fields-]{spaces-over-fields}
\externaldocument[spaces-topologies-]{spaces-topologies}
\externaldocument[spaces-descent-]{spaces-descent}
\externaldocument[spaces-perfect-]{spaces-perfect}
\externaldocument[spaces-more-morphisms-]{spaces-more-morphisms}
\externaldocument[spaces-flat-]{spaces-flat}
\externaldocument[spaces-groupoids-]{spaces-groupoids}
\externaldocument[spaces-more-groupoids-]{spaces-more-groupoids}
\externaldocument[bootstrap-]{bootstrap}
\externaldocument[spaces-pushouts-]{spaces-pushouts}
\externaldocument[spaces-chow-]{spaces-chow}
\externaldocument[groupoids-quotients-]{groupoids-quotients}
\externaldocument[spaces-more-cohomology-]{spaces-more-cohomology}
\externaldocument[spaces-simplicial-]{spaces-simplicial}
\externaldocument[spaces-duality-]{spaces-duality}
\externaldocument[formal-spaces-]{formal-spaces}
\externaldocument[restricted-]{restricted}
\externaldocument[spaces-resolve-]{spaces-resolve}
\externaldocument[formal-defos-]{formal-defos}
\externaldocument[defos-]{defos}
\externaldocument[cotangent-]{cotangent}
\externaldocument[examples-defos-]{examples-defos}
\externaldocument[algebraic-]{algebraic}
\externaldocument[examples-stacks-]{examples-stacks}
\externaldocument[stacks-sheaves-]{stacks-sheaves}
\externaldocument[criteria-]{criteria}
\externaldocument[artin-]{artin}
\externaldocument[quot-]{quot}
\externaldocument[stacks-properties-]{stacks-properties}
\externaldocument[stacks-morphisms-]{stacks-morphisms}
\externaldocument[stacks-limits-]{stacks-limits}
\externaldocument[stacks-cohomology-]{stacks-cohomology}
\externaldocument[stacks-perfect-]{stacks-perfect}
\externaldocument[stacks-introduction-]{stacks-introduction}
\externaldocument[stacks-more-morphisms-]{stacks-more-morphisms}
\externaldocument[stacks-geometry-]{stacks-geometry}
\externaldocument[moduli-]{moduli}
\externaldocument[moduli-curves-]{moduli-curves}
\externaldocument[examples-]{examples}
\externaldocument[exercises-]{exercises}
\externaldocument[guide-]{guide}
\externaldocument[desirables-]{desirables}
\externaldocument[coding-]{coding}
\externaldocument[obsolete-]{obsolete}
\externaldocument[fdl-]{fdl}
\externaldocument[index-]{index}

% Heap Project documents
\externaldocument[linear-algebra-]{linear-algebra}
\externaldocument[groups-]{groups}
\externaldocument[complex-analysis-]{complex-analysis}
\externaldocument[functional-analysis-]{functional-analysis}
\externaldocument[categories-]{categories}
\externaldocument[commutative-algebra-]{commutative-algebra}
\externaldocument[ringed-spaces-]{ringed-spaces}
\externaldocument[homological-algebra-]{homological-algebra}
\externaldocument[lie-algebras-]{lie-algebras}
\externaldocument[representation-theory-]{representation-theory}
\externaldocument[smooth-manifolds-]{smooth-manifolds}
\externaldocument[differential-topology-]{differential-topology}
\externaldocument[algebraic-topology-]{algebraic-topology}
\externaldocument[differential-geometry-]{differential-geometry}
\externaldocument[algebraic-geometry-]{algebraic-geometry}
\externaldocument[symplectic-geometry-]{symplectic-geometry}
\externaldocument[complex-geometry-]{complex-geometry}
\externaldocument[hodge-theory-]{hodge-theory}
\maybeexternaldocument{deformations-}{deformations}
\externaldocument[vertex-algebras-]{vertex-algebras}
\externaldocument[springer-]{springer}
\externaldocument[infty-categories-]{infty-categories}
\externaldocument[seminars-differential-geometry-]{%
seminars-differential-geometry}
\externaldocument[seminars-complex-geometry-]{%
seminars-complex-geometry}
\externaldocument[seminars-algebraic-geometry-]{%
seminars-algebraic-geometry}
\externaldocument[seminars-others-]{%
seminars-others}
\externaldocument[%
geometric-prequantization-]{%
geometric-prequantization}
\externaldocument[surfaces-]{surfaces}
\externaldocument[birational-maps-commutative-algebra-]{birational-maps-commutative-algebra}
\externaldocument[cremona-transformations-]{cremona-transformations}
\externaldocument[derived-categories-of-sheaves-]{derived-categories-of-sheaves}
\externaldocument[geometric-invariant-theory-]{geometric-invariant-theory}
\externaldocument[geometric-stability-conditions-and-group-actions-]{geometric-stability-conditions-and-group-actions}
\externaldocument[moduli-spaces-of-sheaves-]{moduli-spaces-of-sheaves}
\externaldocument[bridgeland-stability-general-theory-]{bridgeland-stability-general-theory}
\externaldocument[hilbert-schemes-]{hilbert-schemes}
\externaldocument[noncommutative-abelian-surfaces-and-kummer-type-hyperkahler-varieties-]{noncommutative-abelian-surfaces-and-kummer-type-hyperkahler-varieties}
\externaldocument[mumford-tate-groups-in-hodge-theory-]{mumford-tate-groups-in-hodge-theory}
\externaldocument[hodge-birational-atoms-2026-]{hodge-birational-atoms-2026}
\externaldocument[xxii-egd-2026-]{%
xxii-egd-2026}
\externaldocument[pde-]{pde}
\externaldocument[probability-]{probability}

\externaldocument[linguistics-]{linguistics}
\externaldocument[genealogia-familiar-]{%
genealogia-familiar}

\externaldocument[%
16th-alga-meeting-2026-]{%
16th-alga-meeting-2026}

\externaldocument[%
iv-jornadas-lefschetz-2026-]{%
iv-jornadas-lefschetz-2026}

\externaldocument[reading-the-masters-]{%
reading-the-masters}

\externaldocument[real-analysis-]{%
real-analysis}

\externaldocument[optimal-transport-]{%
optimal-transport}

\externaldocument[history-]{history}

\externaldocument[seiberg-witten-]{seiberg-witten}

% Theorem environments.
%
\theoremstyle{plain}
\newtheorem{theorem}[subsection]{Theorem}
\newtheorem{proposition}[subsection]{Proposition}
\newtheorem{lemma}[subsection]{Lemma}

\theoremstyle{definition}
\newtheorem{definition}[subsection]{Definition}
\newtheorem{example}[subsection]{Example}
\newtheorem{exercise}[subsection]{Exercise}
\newtheorem{situation}[subsection]{Situation}

\theoremstyle{remark}
\newtheorem{remark}[subsection]{Remark}
\newtheorem{remarks}[subsection]{Remarks}

\numberwithin{equation}{subsection}

% Macros
%
\def\lim{\mathop{\mathrm{lim}}\nolimits}
\def\colim{\mathop{\mathrm{colim}}\nolimits}
\def\Spec{\mathop{\mathrm{Spec}}}
\def\Hom{\mathop{\mathrm{Hom}}\nolimits}
\def\Ext{\mathop{\mathrm{Ext}}\nolimits}
\def\SheafHom{\mathop{\mathcal{H}\!\mathit{om}}\nolimits}
\def\SheafExt{\mathop{\mathcal{E}\!\mathit{xt}}\nolimits}
\def\Sch{\mathit{Sch}}
\def\Mor{\mathop{\mathrm{Mor}}\nolimits}
\def\Ob{\mathop{\mathrm{Ob}}\nolimits}
\def\Sh{\mathop{\mathit{Sh}}\nolimits}
\def\NL{\mathop{N\!L}\nolimits}
\def\CH{\mathop{\mathrm{CH}}\nolimits}
\def\proetale{{pro\text{-}\acute{e}tale}}
\def\etale{{\acute{e}tale}}
\def\QCoh{\mathit{QCoh}}
\def\Ker{\mathop{\mathrm{Ker}}}
\def\Im{\mathop{\mathrm{Im}}}
\def\Coker{\mathop{\mathrm{Coker}}}
\def\Coim{\mathop{\mathrm{Coim}}}

% Boxtimes
%
\DeclareMathSymbol{\boxtimes}{\mathbin}{AMSa}{"02}

%
% Macros for moduli stacks/spaces
%
\def\QCohstack{\mathcal{QC}\!\mathit{oh}}
\def\Cohstack{\mathcal{C}\!\mathit{oh}}
\def\Spacesstack{\mathcal{S}\!\mathit{paces}}
\def\Quotfunctor{\mathrm{Quot}}
\def\Hilbfunctor{\mathrm{Hilb}}
\def\Curvesstack{\mathcal{C}\!\mathit{urves}}
\def\Polarizedstack{\mathcal{P}\!\mathit{olarized}}
\def\Complexesstack{\mathcal{C}\!\mathit{omplexes}}
% \Pic is the operator that assigns to X its picard group, usage \Pic(X)
% \Picardstack_{X/B} denotes the Picard stack of X over B
% \Picardfunctor_{X/B} denotes the Picard functor of X over B
\def\Pic{\mathop{\mathrm{Pic}}\nolimits}
\def\Picardstack{\mathcal{P}\!\mathit{ic}}
\def\Picardfunctor{\mathrm{Pic}}
\def\Deformationcategory{\mathcal{D}\!\mathit{ef}}

\begin{document}
\title{Variations of Hodge Structures}
\maketitle
\tableofcontents

\section{Introduction}
\label{section-introduction}

These notes follow Misha Verbitsky's
2024 course at IMPA.
\href{http://verbit.ru/IMPA/VHS-2024/}
{Course materials} include ten lecture
slide sets, three assignments and an exam.
References below use lecture and slide
numbers. The presentation is a concise
account of the course, supplemented by
period integrals and Picard--Fuchs theory.

A family of varieties gives a local system
of cohomology groups. Its flat connection
records topological transport, while its
Hodge filtration varies holomorphically.
Period maps describe that variation;
monodromy records transport around loops.
The course applies this structure to
Noether--Lefschetz theory.

\section{Hodge structures and polarization}
\label{section-hodge-structures}

Lecture 1, slides 2--5, 7--11, 17--20.
A real Hodge structure of weight $k$ is a
finite-dimensional real vector space with
$$
V_{\mathbb C}=\bigoplus_{p+q=k}V^{p,q},
\qquad
\overline{V^{p,q}}=V^{q,p}.
$$
A rational or integral structure specifies
compatible rational coefficients or a
lattice. The circle acts on $V^{p,q}$
with weight $p-q$; tensor products add
bidegrees.

For a compact K\"ahler manifold, harmonic
representatives give the Hodge decomposition
of cohomology. A polarization is a
nondegenerate rational bilinear pairing,
with parity determined by the weight,
compatible with the decomposition and
satisfying Hodge--Riemann positivity.
For a smooth projective variety, an ample
class and the Lefschetz decomposition
supply polarizations on primitive
cohomology. Orthogonal complements then
show that polarizable rational Hodge
structures form a semisimple category.

Weight-one structures connect the theory
to complex tori. Polarizable integral
structures of type $(1,0)+(0,1)$ correspond
to abelian varieties; passing to rational
coefficients corresponds to working up to
isogeny. Using $H^1$ makes this
correspondence contravariant.

\section{Period domains and period maps}
\label{section-period-maps}

Lecture 2, slides 3--5, 8--18, 25--26.
Use the decreasing filtration convention
$$
F^p=\bigoplus_{a\geq p}V^{a,k-a}.
$$
Fixing Hodge numbers gives a flag variety.
The conjugacy and positivity conditions,
and the equations imposed by a fixed
polarization, define the period domain.
A marking identifies fibre cohomology with
one fixed vector space. The period map
then sends a variety to its Hodge
filtration in this space.

Torelli questions ask how much of the
variety this Hodge data determines.
The lecture discusses tori, curves, K3
surfaces and hyperk\"ahler manifolds.
Markings, polarizations and the relevant
quotients must be specified in each case.
For polarized Hodge structures of K3 type,
the line $V^{2,0}$ determines the filtration.
The period domain consists of lines
satisfying
$$
Q(\omega,\omega)=0,
\qquad Q(\omega,\overline\omega)>0.
$$
When $Q$ has signature $(2,n)$, this is
also the space of oriented positive
real two-planes.

\section{Local systems and Gauss--Manin}
\label{section-gauss-manin}

Lecture 3, slides 2--10.
A local system is a locally constant
sheaf of finite-dimensional vector spaces.
It determines a vector bundle with a flat
connection; conversely, parallel sections
of a flat connection form a local system.
On a connected base, parallel transport
around loops gives its monodromy
representation.

Let $\pi:X\longrightarrow S$ be a smooth
proper submersion. Ehresmann's theorem
identifies nearby fibres differentiably.
Their cohomology groups form the local
system $R^k\pi_*\mathbb C$. Its associated
flat bundle carries the Gauss--Manin
connection.

One construction differentiates a closed
relative form along a lift of a vector
field on the base, then takes its fibrewise
cohomology class. Changing the lift adds
a vertical Lie derivative, which is exact
on closed forms by Cartan's formula.
Thus the induced connection is canonical,
although the chosen lift is not.
Flatness means that transport depends
only on the homotopy class of a path;
it does not require trivial monodromy.

\section{Variations and transversality}
\label{section-variations}

Lecture 3, slides 11--19.
A variation of Hodge structures consists
of a local system with real, rational or
integral structure, and a holomorphic
Hodge filtration on its associated bundle.
Each fibre is a Hodge structure, and the
flat connection satisfies Griffiths
transversality:
$$
\nabla F^p\subset F^{p-1}\otimes\Omega_S^1.
$$
A polarization is parallel and polarizes
each fibre. Smooth projective families
provide the basic examples, using relative
ample classes and primitive cohomology.

The filtration varies holomorphically but
need not be parallel. Transversality says
that differentiation lowers its index by
at most one. On a simply connected base,
flat transport identifies all fibres and
the filtration defines a holomorphic
period map. Otherwise this description
holds on the universal cover, with
monodromy equivariance.

For a complementary treatment, see
\ref{mumford-tate-groups-in-%
hodge-theory-section-variations-of-hs}
in {\it Mumford--Tate groups in Hodge
 theory}.

\section{Fixed parts and semisimplicity}
\label{section-fixed-part}

Lecture 4, slides 4--17;
Lecture 5, slides 6--15.
For a polarized variation over a compact
connected complex base, the fixed part
theorem says that the Hodge components of
a global parallel section are parallel.
Thus monodromy-invariant classes inherit
a Hodge structure. In particular, trivial
monodromy forces the variation to be
constant under these hypotheses.

The proof uses the Chern connection on
the associated graded Hodge bundle,
the components of Gauss--Manin that change
Hodge type, and curvature estimates.
Plurisubharmonic norms become constant on
a compact base. Applied to a Hom variation,
the theorem gives rigidity: a flat map
preserving Hodge structures at one point
preserves them everywhere.

Deligne's semisimplicity theorem, in the
integral polarized setting of Lecture 5,
gives a decomposition over $\mathbb C$:
$$
V=\bigoplus_i L_i\otimes_{\mathbb C}W_i.
$$
Here $L_i$ have irreducible monodromy and
carry complex variations, while $W_i$ are
constant complex Hodge structures.
The decomposition respects Hodge types;
it need not respect the integral lattice.
Lecture 5 also states the extension to
smooth quasiprojective bases, using the
corresponding fixed part theorem.

\section{Hodge groups and K3 type}
\label{section-hodge-groups}

Lecture 6, slides 2--8, 12--15.
The smallest rational algebraic subgroup
containing the Hodge circle is the Hodge
group. The slides call this the
Mumford--Tate group; the convention using
the full Deligne torus also records weight.
Keeping this distinction explicit avoids
confusing the two group conventions.

Tensor invariants recover the Hodge group:
it fixes the rational tensors invariant
under the circle. Hodge substructures
correspond to invariant subspaces.
See \ref{mumford-tate-groups-in-%
hodge-theory-section-mumford-tate-group}
for the group conventions and
\ref{mumford-tate-groups-in-%
hodge-theory-section-tensor-construction}
for tensor constructions.

For a simple polarizable rational Hodge
structure of K3 type, its Hodge
endomorphism algebra is a number field.
Zarhin's theorem makes this field either
totally real or CM; the corresponding
Hodge group is described by orthogonal
or unitary symmetries. The course applies
this to the transcendental Hodge structure,
the smallest rational Hodge substructure
containing the holomorphic top forms.

\section{Noether--Lefschetz loci}
\label{section-noether-lefschetz}

Lecture 3, slide 20;
Lecture 7, slides 6--12;
Lecture 10, slides 4--6, 19.
Fix a rational cohomology class by flat
transport on a simply connected open set.
The condition that it remain of type
$(p,p)$ defines a Hodge locus. In degree
two, these loci detect additional divisor
classes and jumps of the Picard rank.

The Noether--Lefschetz theorem says that a
very general smooth surface of degree
$d\geq4$ in $\mathbb P^3$ has Picard rank
one. Here very general means outside a
countable union of proper algebraic loci.
The hyperplane class supplies that one
class. Irreducibility of monodromy on
primitive cohomology excludes a persistent
proper subsystem of additional divisor
classes. The alternative that all
primitive cohomology has type $(1,1)$ is
excluded by the existence of holomorphic
two-forms for $d\geq4$.

\section{Lefschetz pencils and vanishing
cycles}
\label{section-vanishing-cycles}

Lecture 8, slides 2--11;
Lecture 9, slides 2--18;
Lecture 10, slides 11--19.
A pencil is a one-parameter family of
divisors from a two-dimensional space of
sections. A Lefschetz pencil has smooth
base locus and only ordinary double
points in singular members. Blowing up
the base locus resolves the rational map
to $\mathbb P^1$.

Near a nondegenerate critical point,
the local model is
$$
z_1^2+\cdots+z_{n+1}^2=t.
$$
An $n$-sphere in a smooth fibre contracts
as $t$ approaches zero. Its class is a
vanishing cycle; sweeping it along a path
to zero gives a Lefschetz thimble.
Picard--Lefschetz theory expresses local
monodromy as a rank-one modification
controlled by intersection with this
cycle. The sign depends on dimension
and intersection conventions.

The vanishing cycles generate vanishing
cohomology. For the universal family of
smooth hyperplane sections, monodromy on
this part is irreducible. Applied to the
Veronese embedding, this supplies the
irreducibility used in the proof of the
Noether--Lefschetz theorem.

Picard--Lefschetz theory describes transport
around singular fibres. Picard--Fuchs
equations describe the differential
variation of periods over regular values.
Both concern the same Gauss--Manin local
system from complementary viewpoints.

\section{Periods, Picard--Fuchs equations
and quantum corrections}
\label{section-quantum-periods-background}

\noindent
{\bf Periods and differentiation.}
Over regular values of $f$, follow a
locally constant homology class
$\gamma_t$ in the fibres $C_t=f^{-1}(t)$.
For a relative cohomology class $\omega$,
$$
\Pi(t)=\int_{\gamma_t}\omega_t,
\qquad
\Pi'(t)=\int_{\gamma_t}\nabla_{\partial_t}
\omega.
$$
Gauss--Manin is the connection that
makes this differentiation possible.
Reference: Movasati, Proposition 13.1,
equation (13.6), cited below.
For the local system and Hodge filtration,
see Sections \ref{section-gauss-manin}
and \ref{section-variations}.

\medskip\noindent
{\bf Picard--Fuchs equations.}
In a rank-$r$ cohomology bundle, the
classes $\omega,\nabla\omega,\ldots,
\nabla^r\omega$ are dependent over the
field of meromorphic functions.
Integrating such a relation gives
$$
\sum_{j=0}^{m}b_j(t)\Pi^{(j)}(t)=0,
\qquad m\leq r.
$$
Reference: Movasati, Section 13.5,
equations (13.10)--(13.12).
For compact genus-$g$ curves, $r=2g$.
For affine curves, punctures can add
homology classes, so genus alone does
not give the full rank.
An elliptic curve supplies the familiar
rank-two compact case. The curve is a
fibre of the family; it need not itself
be a moduli space. Periods of a nonzero
holomorphic one-form along a marked
basis of cycles determine its complex
structure through their ratio. This is
the weight-one period map, rather than
an identification of the fibre with the
space parametrizing the family.

For an action integral of a one-form
$\lambda$ on a surface, write locally
$$
d\lambda=df\wedge\eta.
$$
Then differentiation of
$A(t)=\int_{\gamma_t}\lambda$ gives
$$
A'(t)=\int_{\gamma_t}\eta.
$$
This is the Gelfand--Leray derivative.
The rank bound applies directly to the
period of $\eta$. An equation for $A'$
can give an equation for $A$ with one
additional derivative; identifying the
action with a cohomological period
requires the appropriate setup.

\medskip\noindent
{\bf Quantization.}
The relation $[q,p]=h$ replaces commuting
coordinates by noncommuting generators.
The quantum-period construction discussed
in Garay's seminar introduces formal
corrections to the classical area.
For the passage from classical observables
to operators,
see the Dirac rule and Kostant--Souriau
formula in
\ref{geometric-prequantization-%
section-dirac-rule} and
\ref{geometric-prequantization-%
section-kostant-souriau-formula}
in {\it Geometric prequantization}.
For quantum tori and the distinction
between deformation quantization and
quantum cohomology, see
\ref{homological-mirror-symmetry-%
section-where-this-sits} and
\ref{homological-mirror-symmetry-%
section-tori-and-mumford}
in {\it Homological mirror symmetry}.
For the reported quantum-period
construction and its unresolved details,
see \ref{seminars-algebraic-geometry-%
section-garay-quantum-periods}.
These quantum corrections are additional
structure, not part of the definition of
a classical variation of Hodge structures.

\medskip\noindent
{\bf Connections and gauge theory.}
Seiberg--Witten theory studies equations
for connections and spinors modulo gauge
transformations. See
{\it Seiberg--Witten invariants}, Sections
\ref{seiberg-witten-section-dirac} and
\ref{seiberg-witten-section-invariants}.
There the connections are part of the
unknown geometric data. In period theory,
the Gauss--Manin connection describes
variation of fibre cohomology.
The cited sections concern invariants of
four-manifolds; they do not treat
Seiberg--Witten curves and their periods.

\medskip\noindent
{\bf Published reference.}
H. Movasati, {\it A Course in Hodge
Theory: With Emphasis on Multiple
Integrals}, International Press, 2021.
Definition 10.16 (Gelfand--Leray form);
Proposition 13.1 (differentiating periods);
Section 13.5 (Picard--Fuchs equations).
\href{https://w3.impa.br/~hossein/%
myarticles/hodgetheory.pdf}
{Author's text}, consulted
October 14, 2026.

\end{document}
